\scp{Innovations in the lexicon}{08-03}
In the lexicon, there are known contact issues.
\citet{Grimes1991c,Grimes2000a} notes that there are over 12,000 migrants
(15+ generations) from the Sula Islands settled along the north
coast of Buru, and the influence of their language can be seen in some vocabulary
and in the preposed genitive in the vocative of some kin terms
in the Lisela language of north coastal Buru,
but not in the interior or southern dialects of Buru.

Four dialects of Buru average only 27{\%} lexical similarity
with three dialects of Sanana
and Mangole (only trivially higher (2-4{\%}) than Buru shares
with \tsc{Timor-Babar} languages, \tsc{Bima-Lembata}
languages, and even distant Indonesian).
Lisela with Sanana and Mangole shows around 33{\%} similarity,
elevated by significant long-term contact.
Buru with Ambelau is higher at 41--46{\%}.
Buru with Kayeli is trivially higher than that at 44--48{\%},
reflecting the significant contact
and important political role Kayeli has played
on the island for centuries \citep{Grimes2000a}.
Sanana (Fagudu) and Mangole share 72{\%} lexical similarity with
each other \citep{GrimesGrimes1994},
as well as having a few regular differences in their sound correspondences
such as PMP *d > \it{l} in Mangole,
*d > \it{h} in Sanana; the loss of **l between like vowels in
Sanana but not so often in Mangole; PMP *uCu > \it{uCa} in Sanana
but *uCu > \it{oCu} in Mangole (see also \citet{Bloyd2015,Bloyd2020},
and examples in \srf{09-01}).
Taliabu, Soboyo, and Kadai diverge from Buru even further in their
historical phonology, morphology and lexicon.\footnote{
		Grimes’ 220-item word lists collected
		on two trips to the Sula archipelago in 1982
		and 1983 were destroyed in the sectarian violence in Ambon from 1999--2004.
		A perusal of data available from other sources
		\citep{Collins1981,Collins1989,Blust1981,Stokhof1980b}
		does not allow systematic comparison,
		but does indicate that lexical similarity between \tsc{Taliabu}
		and Buru should be significantly lower than the range of 27--33{\%}
		shared between Buru-Lisela and Sanana-Mangole,
		along with many additional differences in
		the historical sound correspondences
		and the morphology.}
So other than sharing *R > \it{h}, it
is misleading by any standard to characterise Buru
and Taliabu as “closely related”
(see \cite[31]{Collins1981}, \citeyear[78, 81, 84]{Collins1982c}, \citeyear[48]{Collins1983}).

The reflexes of *z > \it{y} in \trf{08-08} help reveal patterns
of loans and retentions between \tsc{Taliabu} and Banggai.
\tsc{Taliabu} reflects *z > \it{y},
whereas Banggai reflects *z > \it{d}.
Thus in \trf{08-04}, PMP *qazay ‘chin’ > Banggai \it{ade},
shows that Taliabu \it{n-ade} is a loan from Banggai.

Like Balantak, Soboyo-Taliabu-Kadai \it{tuma} ‘headlouse’ < PMP
*tumah ‘clothes louse, body louse’.
In contrast, \tsc{Sula-Buru} languages reflect *kutu ‘headlouse’.
Like \tsc{Saluan-Banggai} languages, Taliabu takes its word for
‘fly (insect)’ from *laləj,
whereas all \tsc{Sula-Buru} languages take their word for ‘fly
(insect)’ from \tsc{PAn} *bəRŋaw ‘bluebottle,
large fly’ (see \trf{08-06})

Soboyo, Taliabu, and Kadai have developed the innovations in \qf{08-01}
in the lexicon, distinct from Banggai and Balantak,
and also from \tsc{Sula-Buru} languages.

{\wg\st{6pt}\begin{xltabular}[l]{\linewidth}{>{\hspace{-\tabcolsep}}l<{\hspace{-\tabcolsep}}>{\hspace{-\tabcolsep}\enspace}llllX}
\lsptoprule
\tbx{08-01}&	{Soboyo}	&	{Taliabu}	&	{Kadai}	&	{gloss}	&	{note}	\\ \midrule
\endfirsthead
\lsptoprule
					&	{Soboyo}	&	{Taliabu}	&	{Kadai}	&	{gloss}	&	{note}	\\ \midrule
\endhead
&	howo	&	howo	&	howo	&	two	&	< PMP *duha (see \trf{08-12})	\\
&	dina	&	dina	&	dina	&	sun	&	\hp{\,}	\\
&	n-dolu	&	n-dolu	&	n-dolu	&	stone	&	\hp{\,}	\\
&	utoin	&	utoɲ	&	uto	&	fire	&	\hp{\,}	\\
&	tala-in	&	talaŋ	&	n-tala	&	bone	&	(see \trf{08-13})	\\
&	mboa	&	mendadi	&	mboa	&	pig	&	\it{mboa} possibly irr. from *bəRək `domesticated pig'	\\
&	mine	&	mini	&	mine	&	\mc{2}{l}{hear, listen}	\\
&	ŋ-gomoɲ	&	naŋ-gomo	&	ŋ-gomo	&	cold	&	\hp{\,}	\\
&	barede	&	beridi	&	barede	&	black	&	\hp{\,}	\\
&	wohe	&	wohi	&	wohe	&	wind	&	\hp{\,}	\\
&	poloi-t	&	polo	&	polo	&	blood	&	\hp{\,}	\\
&	kena	&	ikan	&	kena	&	fish	&	Mangole and Sanana \it{kena} also `fish'	\\
&	lage	&	lagi	&	lage	&	boat	&	Ternate loan	\\
&	salak sia	&	ʧala sia	&	salak sia	&	\mc{2}{l}{one thousand}	\\
\lspbottomrule
\end{xltabular}\addtocounter{table}{-1}}

Soboyo-Taliabu-Kadai all share \it{bira} ‘rice’,
ultimately from PMP *bəRas ‘husked rice’,
but probably via an intermediate source.\footnote{
		The regular
		sound changes in \tsc{Taliabu}
		would be *b > \it{f},
		*ə > \it{o}, *R > \it{h}.
		Many languages in Sulawesi have /i/ in the first syllable. Cf.
		Proto-Sangiric **biRas ‘rice’, Selayar [ˈbɪːrasa] ‘uncooked rice’
		\citep[106]{GrimesGrimes1987}, and others.}
However \tsc{Sula-Buru} languages have forms from
PMP *pajay via regular sound correspondences: Buru \it{pala}
‘rice’, Ambelau \it{fala} ‘rice’, and \tsc{Ambon-Seram} Kayeli \it{hala} ‘rice’.
Lisela \it{hala} ‘rice’ is a Kayeli loan.
Sanana \it{bira} is a \tsc{Taliabu} loan.
Mangole has \it{pamasi} ‘rice’ from an unknown source.

All languages in both \tsc{Taliabu}
and \tsc{Sula-Buru} get their words for ‘banana’ from a form
like *biak: Soboyo \it{fiaʔ},
Taliabu, Kadai \it{fia}. (see \trf{07-13}).
In the lexicon, all \tsc{Sula-Buru} languages get their word
for ‘eye, source’ from Proto-Sula-Buru **dama (see \trf{07-13}).
None of the \tsc{Taliabu} languages share that with them,
with Soboyo \it{bele-ŋ} ‘eye’,
Taliabu (Likitobi) \it{bele-n} ‘eye’,
Kadai \it{nata} ‘eye’ (possibly from /n-mata/).
\tsc{Sula-Buru} **kau ‘tree, wood’ distinguishes it from \tsc{Taliabu}
**kayu, since Mangole-Sanana normally preserve *y,
as shown in \trf{07-13}.\footnote{
		Sanana, which normally preserves *y > \it{y}, has \it{kau}.
		\citet[77]{Mead2003b} proposes **kayu for Proto-Saluan-Banggai,
		also supported by the fortition of *y in **kayu in reflexes in
		nearby \tsc{Bungku-Tolaki} languages, but notes the
		possibility of an alternate view in the (unlikely)
		event that Banggai \it{kayuŋ} were to be a Malay loan.}



